\documentclass[../../main.tex]{subfiles}
\begin{document}

\begin{flushright}
\footnotesize\textit{【自動文字起こし・要確認】}
\end{flushright}

{\bf [解]}
$f(x)=\dfrac{ax+b}{x^2+x+1}$とする．$y=f(x)$とおけば，与式は
\begin{align}
y^3-2y^2-y+2\ge0
\end{align}
すなわち
\begin{align}
(y-1)(y^2-y-2)\ge0 \iff (y-1)(y-2)(y+1)\ge0
\end{align}
\begin{align}
\therefore\ -1\le y\le1\ \text{または}\ 2\le y\quad\cdots①
\end{align}

したがって，全ての$x$で①が成立すれば良い．$f(x)\xrightarrow{x\to\pm\infty}0$かつ，$x^2+x+1>0$から$f$が連続であることより，$2\le y$となる$x$がある時，①が全ての$x$で成立することはなく適さない．

よって，$-1\le y\le1$のみ調べる．
\begin{align}
-(x^2+x+1)\le ax+b\le x^2+x+1
\end{align}
\begin{align}
-x^2-(a+1)x-1\le b\le x^2+(1-a)x+1\quad\cdots②
\end{align}

②の右辺の最小値は$x=\dfrac{a-1}{2}$の時の$1-\left(\dfrac{1-a}{2}\right)^2=\dfrac{-a^2+2a+3}{4}$，
左辺の最大値は$x=-\dfrac{a+1}{2}$の時の$\dfrac{(a+1)^2}{4}-1=\dfrac{a^2+2a-3}{4}$だから，全ての$x$で②が成立する条件は，
\begin{align}
\dfrac{a^2+2a-3}{4}\le b\le\dfrac{-a^2+2a+3}{4}．
\end{align}

これを図示して，以下斜線部（境界含む）．

\begin{figure}[htb]
\centering
\begin{tikzpicture}[scale=1.15, >=stealth]
  \begin{axis}[
      axis lines=middle,
      xlabel={$a$}, ylabel={$b$},
      xmin=-2.3, xmax=2.3, ymin=-1.6, ymax=1.6,
      xtick={-1.7320508,1.7320508}, xticklabels={$-\sqrt{3}$,$\sqrt{3}$},
      ytick=\empty,
      clip=false,
      domain=-1.7320508:1.7320508, samples=100, smooth,
    ]
    \addplot[name path=upper, thick] {(-x^2+2*x+3)/4};
    \addplot[name path=lower, thick] {(x^2+2*x-3)/4};
    \addplot[gray!25] fill between[of=upper and lower];
  \end{axis}
\end{tikzpicture}
\caption{条件②を満たす$(a,b)$の範囲（斜線部）}
\end{figure}

\newpage

\end{document}
